\begin{table}[ht]
\centering
\caption{Sensitivity of the annual procurement-cost implication to the admissibility calibration.}
\label{tab:procurement_cost_sensitivity}
\begin{threeparttable}
\small
\begin{tabular}{lccc}
\toprule
Admissible share & Midpoint & Lee lower bound & Lee upper bound \\
\midrule
25\% & \$13.9 M / yr & \$11.9 M / yr & \$15.8 M / yr \\
50\% & \$27.8 M / yr & \$23.9 M / yr & \$31.7 M / yr \\
75\% & \$41.6 M / yr & \$35.8 M / yr & \$47.5 M / yr \\
\bottomrule
\end{tabular}
\begin{tablenotes}[flushleft]\footnotesize
\item \textit{Notes:} Each row applies the bounded litigated-over-administrative price gap to annual litigated spending under the listed admissible share. The middle row is the calibration used in Table~\ref{tab:procurement_cost_bound}. The table is a fiscal procurement-cost sensitivity calculation, not a full welfare estimate.
\end{tablenotes}
\end{threeparttable}
\end{table}
